\section{Statistical analysis}
\subsection{Exploration and model freezing}
Use polychoric correlations, ordinal parallel analysis based on independent within-item permutations preserving marginal frequencies (500 replicates, 95th percentile), minimum-residual extraction and oblique rotation. Compare plausible neighbouring factor counts, report pattern and structure matrices, communalities and factor correlations. Flag primary loading below .40, secondary loading at least .30 with a primary-secondary difference below .20, sparse categories and local dependence. Content coverage and interpretability govern decisions alongside statistical criteria; do not optimise reliability by deleting unique content. Retain at least three content-distinct indicators per domain; target four or five. If no stable interpretable solution emerges, revise and recruit new development data. No CFA data are used for item selection.

\subsection{Confirmation and reliability}
For item $j$ and person $i$, use an ordinal probit model
\[
Y_{ij}=k \iff \tau_{j,k}<Y^{*}_{ij}\leq\tau_{j,k+1},\qquad
Y^{*}_{ij}=\boldsymbol{\lambda}_{j}^{\mathsf T}\boldsymbol{\eta}_i+\epsilon_{ij}.
\]
Order thresholds, identify latent variances and use a robust diagonally weighted least-squares estimator with mean/variance-adjusted inference \parencite{flora2004}. Compare frozen correlated and one-factor models using fit indices, residuals and model plausibility, not ordinary subtraction of robust chi-square statistics. Report robust chi-square, degrees of freedom, CFI, TLI, RMSEA with 90\% interval, SRMR, standardised loadings and factor correlations with intervals. CFI/TLI near .95, RMSEA near .06 and SRMR near .08 are descriptive reference values, not universal cutoffs. Improper estimates, weak identification and near-unit factor correlations require explicit interpretation. No modification-index-driven repairs count as confirmatory evidence.

For a simple congeneric domain report model-based omega on the latent response scale, $\omega=(\mathbf{1}^{\mathsf T}\lambda)^2/[(\mathbf{1}^{\mathsf T}\lambda)^2+\mathbf{1}^{\mathsf T}\Theta\mathbf{1}]$ with unit factor variance. Distinguish it from observed ordinal sum-score reliability; estimate the latter using model-implied category probabilities if feasible. Bootstrap 2000 participant resamples for intervals, reporting failures. Report same-episode test--retest absolute agreement ICC(A,1) with confidence intervals and paired differences.

\subsection{Validity and invariance}
Estimate H3--H5 using frozen domain scores, Spearman correlations, partial rank associations for clarity, and participant bootstrap intervals (5000 resamples). Holm adjustment across six planned validity contrasts (five associations and one paired correlation difference); report estimates and unadjusted intervals. If a comparator is unavailable, state that explicitly and retain its place in the family with a non-rejection rather than replacing it post hoc. Assess discriminant evidence via factor correlations and relationships to distinct constructs. Test ordinal multi-group invariance across prespecified musical training categories (at least one year formal training versus less/no training) only with adequate group sizes and category support established by simulation. Follow identification appropriate to ordinal thresholds and loadings \parencite{wu2016}. Do not compare latent means without sufficient invariance.

\subsection{Missingness and sample size}
Report eligibility, nonresponse and missingness by item, domain and recruitment source. Treat missing as missing, never zero; no forced completion. Prespecify complete-item CFA as a transparent baseline with selection limitations; evaluate ordinal-compatible multiple imputation under MAR and pattern-specific MNAR sensitivity (shifts of latent thresholds or latent imputations) before locking the analysis. Record estimator implementation and simulation performance; do not claim calibration yet. Domain means may be used for external validity only after unidimensionality and scoring justification; require at least 80\% of retained domain items. No total score or diagnostic thresholds.

The EFA and CFA targets are provisional. Before launch simulate ordinal models at multiple sample sizes varying loadings, factor correlations, category imbalance and missingness \parencite{muthen2002}. Evaluate recovery, convergence, improper solutions, coverage, bias, power for validity hypotheses and group-size sensitivity. Require at least 95\% proper converged fits, absolute loading bias at most .05, approximate 95\% interval coverage of .92--.98 and at least .80 power for prespecified meaningful effects in core scenarios. Stress scenarios and failures are reported. Do not reduce required sample size after looking at real effect estimates. Fisher-transform correlation planning is not factor-model power.


\subsection{Computational development and release gates}
Development code implements polychoric likelihoods, ordinal exploratory analysis, threshold-aware diagonal weighting with sandwich covariance, reliability, selection sensitivity, ordinal probit chained imputation and the six external validity contrasts. A conditional-normal quadrature reference checks polychoric probabilities independently of the primary Plackett integration formulation. The custom moment-matched test and discrepancy indices are not validated WLSMV substitutes. Affine-linked group diagnostics are not a full nested multigroup CFA. Imputed-item trace diagnostics do not establish joint posterior convergence. Before confirmatory freeze, validate the intended estimator against an independent implementation, complete the prospective simulation grid and obtain independent scientific review. Report synthetic calibration separately from empirical evidence. No empirical scoring authorization follows from passing software tests.
